\paragrafos{Ιδιότητες ισοσκελούς και ισόπλευρου τριγώνου}{Β.3.5, «Ιδιότητες ισοσκελούς και ισόπλευρου τριγώνου»}

\begin{anagke}
Η Αγγελική σχεδιάζει στο τετράδιό της ένα έλατο με τρίγωνα: κάθε τρίγωνο έχει δύο ίσες πλευρές. Τα κόβει με το ψαλίδι και τα διπλώνει στη μέση, και παρατηρεί ότι οι δύο μισές πέφτουν ακριβώς η μία πάνω στην άλλη. «Και οι γωνίες στη βάση είναι ίσες;», ρωτά η Ζωή. «Γιατί να είναι;» Η Πολυξένη το γράφει στο τετράδιό της με κωδικούς, γιατί της αρέσουν τα μοτίβα: ισοσκελές $\to$ ίσες γωνίες βάσης. Ο Ηλίας το μετρά σε τρία διαφορετικά ισοσκελή: τρεις φορές βγαίνουν ίσες.

Η κυρία Αλεξάνδρα ρωτά: «Η πτυχή του χαρτιού είναι και ύψος του τριγώνου; Και διχοτόμος; Και διάμεσος;». Η παρέα δοκιμάζει με ισόπλευρο τρίγωνο και βρίσκει ότι δεν υπάρχει μία πτυχή, αλλά τρεις.

Σε αυτή την παράγραφο μαθαίνουμε γιατί το ισοσκελές τρίγωνο έχει άξονα συμμετρίας, ποιες γωνίες του είναι ίσες, και ότι σε αυτό η διάμεσος, το ύψος και η διχοτόμος προς τη βάση είναι το ίδιο τμήμα. Θα δούμε και τις αντίστοιχες ιδιότητες του ισόπλευρου τριγώνου.
\end{anagke}

\begin{lexi}{βάση, προσκείμενες γωνίες, ταυτίζονται}{ονόματα και ιδιότητες}
Η \lexiword{βάση} του ισοσκελούς είναι η πλευρά που δεν είναι από τις δύο ίσες· στηρίζεται πάνω της σαν στο έδαφος. \lexiword{Προσκείμενες} (προς + κείμαι) λέμε τις γωνίες που «κείτονται πάνω» στην ίδια πλευρά, στα άκρα της. \lexiword{Ταυτίζονται} σημαίνει ότι δύο πράγματα είναι ένα και το ίδιο: δύο τμήματα που ταυτίζονται δεν είναι δύο, είναι ένα.
\end{lexi}

\begin{metafrasi}
Το ισοσκελές τρίγωνο με ίσες πλευρές $\mathrm{ΑΒ}$ και $\mathrm{ΑΓ}$. & $\mathrm{ΑΒ}=\mathrm{ΑΓ}$, βάση η $\mathrm{ΒΓ}$\\
Οι γωνίες στη βάση είναι ίσες. & $\hat{\mathrm{Β}}=\hat{\Gamma}$\\
Η διάμεσος $\mathrm{ΑΜ}$ προς τη βάση είναι και ύψος. & $\mathrm{ΒΜ}=\mathrm{ΜΓ}$ και $\mathrm{ΑΜ}\perp\mathrm{ΒΓ}$\\
Η διάμεσος $\mathrm{ΑΜ}$ είναι και διχοτόμος. & $\hat{\mathrm{Α}}_1=\hat{\mathrm{Α}}_2$\\
Το ισόπλευρο τρίγωνο. & $\mathrm{ΑΒ}=\mathrm{ΒΓ}=\mathrm{ΓΑ}$\\
\end{metafrasi}

\tmima{Η θεωρία}

\tmimasmall{Ισοσκελές τρίγωνο}

Ένα ισοσκελές τρίγωνο $\mathrm{ΑΒΓ}$ έχει $\mathrm{ΑΒ}=\mathrm{ΑΓ}$. Το $\mathrm{Α}$ είναι η \textbf{κορυφή} και η $\mathrm{ΒΓ}$ η \textbf{βάση} του. Αν πάρουμε το μέσο $\mathrm{Μ}$ της βάσης και φέρουμε την $\mathrm{ΑΜ}$, διπλώνοντας το σχήμα κατά την $\mathrm{ΑΜ}$ το $\mathrm{Β}$ πέφτει πάνω στο $\Gamma$.

\begin{center}
\begin{tikzpicture}[scale=0.95]
\draw[thick] (0,0)--(4,0)--(2,3)--cycle;
\draw[gblue,thick,dashed] (2,3.4)--(2,-0.4);
\node[below left] at (0,0) {$\mathrm{Β}$}; \node[below right] at (4,0) {$\Gamma$}; \node[above right] at (2,3) {$\mathrm{Α}$}; \node[below right] at (2,0) {$\mathrm{Μ}$};
\draw (2,0.3)--(2.3,0.3)--(2.3,0);
\draw (0.9,1.5)--(1.1,1.4); \draw (2.9,1.4)--(3.1,1.5); \draw (1,-0.1)--(1,0.1); \draw (3,-0.1)--(3,0.1);
\draw (0.6,0) arc(0:56:0.6); \draw (3.4,0) arc(180:124:0.6);
\end{tikzpicture}
\end{center}

\begin{kanonas}{ιδιότητες ισοσκελούς τριγώνου}
Σε ισοσκελές τρίγωνο $\mathrm{ΑΒΓ}$ με $\mathrm{ΑΒ}=\mathrm{ΑΓ}$:
\begin{itemize}[leftmargin=1.2em,itemsep=0.3ex,topsep=2pt]
\item Η ευθεία της διαμέσου $\mathrm{ΑΜ}$ προς τη βάση είναι \textbf{άξονας συμμετρίας} του τριγώνου.
\item Οι προσκείμενες στη βάση γωνίες είναι \textbf{ίσες}: $\hat{\mathrm{Β}}=\hat{\Gamma}$.
\item Η διάμεσος, το ύψος και η διχοτόμος προς τη βάση \textbf{ταυτίζονται}: είναι το ίδιο τμήμα $\mathrm{ΑΜ}$.
\end{itemize}
\end{kanonas}

Γιατί ισχύουν. Στη συμμετρία ως προς τον άξονα $\mathrm{ΑΜ}$ το $\mathrm{Β}$ γίνεται $\Gamma$, και το $\mathrm{Α}$ μένει στη θέση του. Έτσι η γωνία $\hat{\mathrm{Β}}$ πέφτει πάνω στη $\hat{\Gamma}$ και είναι ίσες. Η $\hat{\mathrm{Α}}_1$ πέφτει πάνω στην $\hat{\mathrm{Α}}_2$, άρα η $\mathrm{ΑΜ}$ είναι διχοτόμος. Και η γωνία $\mathrm{ΑΜ}\hat{\mathrm{Β}}$ πέφτει πάνω στην $\mathrm{ΑΜ}\hat{\Gamma}$, άρα είναι ίσες και μαζί κάνουν ευθεία γωνία: καθεμία είναι $90^\circ$, δηλαδή η $\mathrm{ΑΜ}$ είναι και ύψος.

Η ιδιότητα αυτή συνδέεται με τη μεσοκάθετο: αφού $\mathrm{ΑΒ}=\mathrm{ΑΓ}$, το $\mathrm{Α}$ ισαπέχει από τα $\mathrm{Β}$ και $\Gamma$, άρα βρίσκεται στη μεσοκάθετο της $\mathrm{ΒΓ}$. Και η $\mathrm{ΑΜ}$ είναι ακριβώς αυτή η μεσοκάθετος.

\tmimasmall{Ισόπλευρο τρίγωνο}

Το ισόπλευρο τρίγωνο έχει και τις τρεις πλευρές του ίσες. Αφού κάθε πλευρά του μπορεί να παίξει τον ρόλο της βάσης, όσα ισχύουν για το ισοσκελές ισχύουν και για καθεμία από τις τρεις πλευρές.

\begin{center}
\begin{tikzpicture}[scale=0.95]
\draw[thick] (0,0)--(4,0)--(2,3.464)--cycle;
\draw[gblue,dashed] (2,3.9)--(2,-0.4); \draw[gblue,dashed] (4,0)--(0.5,2.598)--(-0.3,3.17); \draw[gblue,dashed] (0,0)--(3.5,2.598)--(4.3,3.17);
\node[below left] at (0,0) {$\mathrm{Β}$}; \node[below right] at (4,0) {$\Gamma$}; \node[above] at (2,3.9) {$\mathrm{Α}$};
\end{tikzpicture}
\end{center}

\begin{kanonas}{ιδιότητες ισόπλευρου τριγώνου}
\begin{itemize}[leftmargin=1.2em,itemsep=0.3ex,topsep=2pt]
\item Οι ευθείες των τριών διαμέσων είναι \textbf{άξονες συμμετρίας} του τριγώνου: έχει τρεις άξονες συμμετρίας.
\item Όλες οι γωνίες του είναι \textbf{ίσες}. Θα αποδείξουμε στην 3.9 ότι κάθε μία είναι $60^\circ$.
\item Η διάμεσος, το ύψος και η διχοτόμος που αντιστοιχούν σε οποιαδήποτε πλευρά \textbf{ταυτίζονται}.
\end{itemize}
\end{kanonas}

\begin{prosoxi}
\begin{itemize}[leftmargin=1.2em,itemsep=0.3ex,topsep=2pt]
\item Οι ίσες γωνίες του ισοσκελούς είναι οι δύο γωνίες της \emph{βάσης}. Η γωνία της κορυφής είναι γενικά διαφορετική από αυτές.
\item Στο ισοσκελές η διάμεσος προς τη βάση ταυτίζεται με το ύψος και τη διχοτόμο. Η διάμεσος προς μία από τις ίσες πλευρές δεν είναι κατά κανόνα ύψος ούτε διχοτόμος.
\item Το ισοσκελές τρίγωνο που δεν είναι ισόπλευρο έχει έναν μόνο άξονα συμμετρίας, ενώ το ισόπλευρο έχει τρεις.
\item Η ιδιότητα $\hat{\mathrm{Β}}=\hat{\Gamma}$ ισχύει όταν $\mathrm{ΑΒ}=\mathrm{ΑΓ}$. Σε ένα σκαληνό τρίγωνο οι γωνίες δεν είναι ίσες, όσο κι αν στο σχέδιο μοιάζουν.
\end{itemize}
\end{prosoxi}

\tmima{Λυμένα παραδείγματα}

\paradeigma{Οι γωνίες της βάσης}
\tixero Το τρίγωνο $\mathrm{ΑΒΓ}$ είναι ισοσκελές με $\mathrm{ΑΒ}=\mathrm{ΑΓ}$ και $\hat{\mathrm{Β}}=70^\circ$.\par
\tithelo Η γωνία $\hat{\Gamma}$.\par
\possindeo Οι προσκείμενες στη βάση $\mathrm{ΒΓ}$ γωνίες του ισοσκελούς είναι ίσες.\par
\apantisi $\hat{\Gamma}=\hat{\mathrm{Β}}=70^\circ$.\par
\elegxos Σε σχέδιο με μοιρογνωμόνιο μετράμε και τις δύο γωνίες της βάσης: βγαίνουν $70^\circ$.

\paradeigma{Η διάμεσος του ισοσκελούς}
\tixero Το τρίγωνο $\mathrm{ΑΒΓ}$ είναι ισοσκελές με $\mathrm{ΑΒ}=\mathrm{ΑΓ}$. Η $\mathrm{ΑΜ}$ είναι διάμεσος, $\mathrm{ΒΓ}=10$ cm και $\hat{\mathrm{Α}}=50^\circ$.\par
\tithelo Τα $\mathrm{ΒΜ}$, $\mathrm{ΑΜ}\hat{\mathrm{Β}}$ και $\mathrm{ΒΑΜ}$.\par
\possindeo Η διάμεσος προς τη βάση είναι και ύψος και διχοτόμος. Το $\mathrm{Μ}$ είναι μέσο της $\mathrm{ΒΓ}$.\par
\apantisi $\mathrm{ΒΜ}=10:2=5$ cm. Είναι ύψος: $\mathrm{ΑΜ}\hat{\mathrm{Β}}=90^\circ$. Είναι διχοτόμος: $\mathrm{ΒΑΜ}=50^\circ:2=25^\circ$.\par
\elegxos Οι δύο μισές γωνίες $25^\circ+25^\circ=50^\circ$.

\paradeigma{Οι άξονες}
\tixero Ένα ισόπλευρο τρίγωνο πλευράς $6$ cm.\par
\tithelo Πόσους άξονες συμμετρίας έχει και ποιοι είναι;\par
\possindeo Κάθε διάμεσος είναι άξονας συμμετρίας και έχει τρεις διαμέσους.\par
\apantisi Τρεις άξονες: οι ευθείες των τριών διαμέσων, που περνούν από το ίδιο σημείο.\par
\elegxos Αν διπλώσω το τρίγωνο κατά οποιονδήποτε από τους τρεις άξονες, οι δύο μισές συμπίπτουν.

\paradeigma{Ποιο τρίγωνο διπλώνεται;}
\tixero Η Ρήνη διπλώνει ένα χάρτινο τρίγωνο ώστε το $\mathrm{Β}$ να πέσει πάνω στο $\Gamma$, και οι δύο μισές συμπίπτουν ακριβώς.\par
\tithelo Τι είδους τρίγωνο είναι;\par
\possindeo Η πτυχή είναι άξονας συμμετρίας στον οποίο το $\mathrm{Β}$ γίνεται $\Gamma$ και το $\mathrm{Α}$ μένει στη θέση του, άρα $\mathrm{ΑΒ}=\mathrm{ΑΓ}$.\par
\apantisi Ισοσκελές, με κορυφή το $\mathrm{Α}$.\par
\elegxos Η πτυχή περνά από το $\mathrm{Α}$ και από το μέσο της $\mathrm{ΒΓ}$.

\tmima{Ασκήσεις}

\zoni{Το κατάλαβα}{ισοσκελές και ισόπλευρο}

\ex Σωστό ή λάθος; α) Οι προσκείμενες στη βάση γωνίες ισοσκελούς τριγώνου είναι ίσες. \quad β) Σε κάθε ισοσκελές η διάμεσος προς τη βάση είναι και διχοτόμος. \quad γ) Σε κάθε ισόπλευρο τρίγωνο οι ευθείες των υψών είναι και μεσοκάθετοι των πλευρών του. \quad δ) Το ισοσκελές τρίγωνο έχει πάντα τρεις άξονες συμμετρίας. \quad ε) Σε κάθε ισοσκελές η διάμεσος προς τη βάση είναι άξονας συμμετρίας.

\ex Το τρίγωνο $\mathrm{ΑΒΓ}$ είναι ισοσκελές με $\mathrm{ΑΒ}=\mathrm{ΑΓ}$ και η $\mathrm{ΑΜ}$ είναι ύψος. Αν $\mathrm{ΒΓ}=8$ cm και $\hat{\mathrm{Α}}=50^\circ$, βρες τα $\mathrm{ΒΜ}$, $\mathrm{ΑΜ}\hat{\mathrm{Β}}$ και $\mathrm{ΒΑΜ}$.

\ex Στο ισοσκελές τρίγωνο $\mathrm{ΑΒΓ}$ με κορυφή το $\mathrm{Α}$ δίνονται: α) $\hat{\mathrm{Β}}=70^\circ$, β) $\hat{\Gamma}=48^\circ$. Βρες σε κάθε περίπτωση την άλλη γωνία της βάσης και δικαιολόγησε.

\ex Στο ισόπλευρο τρίγωνο $\mathrm{ΑΒΓ}$ η $\mathrm{ΑΔ}$ είναι διάμεσος και $\mathrm{ΑΒ}=6$ cm. Βρες το $\mathrm{ΒΓ}$, το $\mathrm{ΒΔ}$ και τη γωνία $\mathrm{ΑΔ}\hat{\mathrm{Β}}$. Πόσο είναι η γωνία $\mathrm{ΒΑΔ}$ αν ξέρουμε ότι $\hat{\mathrm{Α}}=60^\circ$;

\ex Κατασκεύασε με κανόνα και διαβήτη ένα ισοσκελές τρίγωνο με $\mathrm{ΑΒ}=\mathrm{ΑΓ}=5$ cm και $\mathrm{ΒΓ}=6$ cm. Μέτρησε τις γωνίες της βάσης και βρες τον άξονα συμμετρίας με δίπλωση.

\ex Κατασκεύασε με κανόνα και διαβήτη ένα ισόπλευρο τρίγωνο με πλευρά $5$ cm. Φέρε τους τρεις άξονες συμμετρίας και μέτρησε τις γωνίες του.

\ex Πόσους άξονες συμμετρίας έχει: α) ένα ισοσκελές τρίγωνο που δεν είναι ισόπλευρο, β) ένα ισόπλευρο, γ) ένα σκαληνό;

\zoni{Το χρησιμοποιώ}{στέγες, πανιά, πινακίδες}

\ex \textbf{Η στέγη.} Η στέγη του σπιτιού της γιαγιάς Μυρτώς έχει σχήμα ισοσκελούς τριγώνου με γωνία βάσης $35^\circ$. Πόση είναι η άλλη γωνία της βάσης;

\ex \textbf{Το πανί.} Ένα ισοσκελές τριγωνικό πανί έχει βάση $12$ cm. Η διάμεσος από την κορυφή προς τη βάση έχει μήκος $8$ cm. Πόσο απέχει το πόδι της από κάθε άκρο της βάσης; Τι γωνία σχηματίζει με τη βάση;

\ex \textbf{Η σημαία.} Στο ισοσκελές τρίγωνο $\mathrm{ΑΒΓ}$ με $\mathrm{ΑΒ}=\mathrm{ΑΓ}$ η $\mathrm{ΑΔ}$ είναι η διχοτόμος της γωνίας $\hat{\mathrm{Α}}=64^\circ$. Βρες τη γωνία $\mathrm{ΒΑΔ}$ και τη γωνία $\mathrm{ΑΔ}\hat{\mathrm{Β}}$.

\ex \textbf{Η πινακίδα.} Η πινακίδα «Δώσε προτεραιότητα» έχει σχήμα ισόπλευρου τριγώνου πλευράς $40$ cm. Βρες την περίμετρο και πόσους άξονες συμμετρίας έχει.

\ex \textbf{Η δίπλωση του Στέλιου.} Ο Στέλιος διπλώνει ένα τρίγωνο και οι δύο μισές συμπίπτουν. Λέει: «Άρα το τρίγωνο είναι ισόπλευρο!». Έχει δίκιο; Τι ακριβώς είναι;

\zoni{Το σκέφτομαι}{γιατί συμβαίνει}

\ex Εξήγησε γιατί οι γωνίες στη βάση ενός ισοσκελούς είναι ίσες, με βάση τη δίπλωση.

\ex Γιατί η ευθεία $\mathrm{ΑΜ}$ στο ισοσκελές τρίγωνο είναι η μεσοκάθετος της βάσης $\mathrm{ΒΓ}$; (Σκέψου το $\mathrm{Α}$ και τις αποστάσεις του από τα $\mathrm{Β}$, $\Gamma$.)

\ex Σε ισοσκελές τρίγωνο $\mathrm{ΑΒΓ}$ ($\mathrm{ΑΒ}=\mathrm{ΑΓ}$) φέρε τη διάμεσο προς την πλευρά $\mathrm{ΑΓ}$. Είναι και ύψος; Είναι και διχοτόμος; Έλεγξε με σχέδιο.

\ex Σχεδίασε ένα ορθογώνιο ισοσκελές τρίγωνο με κάθετες πλευρές $4$ cm. Μέτρησε τις δύο άλλες γωνίες. Τι παρατηρείς;

\zoni{Για το σπίτι}{με χαρτί}

\ex Κόψε ένα ισοσκελές τρίγωνο από χαρτί, δίπλωσέ το στη μέση και μέτρησε τις γωνίες της βάσης.

\ex Κόψε ένα ισόπλευρο τρίγωνο από χαρτί και βρες με δίπλωση τους τρεις άξονες συμμετρίας του. Πού τέμνονται;
